\subsection{Quotient of a groupoid by the action of a group}

Let G be a transitive groupoid, let K be a group, and let \(\theta\colon\mathrm{K}\to\mathrm{Aut}(\mathrm{G})^{\circ}\) be a group homomorphism from K into the opposite group of the automorphism group of the groupoid G. We will say that the group K acts on the right on G. If \(k\in\mathrm{K}\) we will sometimes denote by \(k^{*}x\) (resp. \(k^{*}f\) ) the image of a vertex \(x\) (resp. of an arrow \(f\) ) of G under the groupoid automorphism \(\theta(k)\) .

Let \(|\mathrm{K}|\) be the groupoid whose set of vertices is \(\mathrm{K}\) and whose set of arrows connecting two vertices is empty if these vertices are distinct, and is reduced to a single element otherwise. Let \(\mathrm{H}\) be the product groupoid \(\mathrm{G} \times |\mathrm{K}|\) ; a vertex of \(\mathrm{H}\) is a pair \((a,k)\) , where \(a\) is a vertex of \(\mathrm{G}\) and \(k\) is an element of \(\mathrm{K}\) ; if \(f\) is an arrow of \(\mathrm{G}\) joining a vertex \(a\) to a vertex \(b\) , we shall denote \((f,k)\) the unique arrow of \(\mathrm{H}\) connecting \((a,k)\) to \((b,k)\) . Let \(\varphi: \mathrm{H} \to \mathrm{G}\) be the groupoid morphism given by the first projection, and let \(\psi: \mathrm{H} \to \mathrm{G}\) be the groupoid morphism such that \(\operatorname{Som}(\psi)((a,k)) = k^*a\) and \(\operatorname{Fl}(\psi)((f,k)) = k^*f\) if \(k \in \mathrm{K}\) , \(a \in \operatorname{Som}(\mathrm{G})\) and \(f \in \operatorname{Fl}(\mathrm{G})\) .

Let G/K denote the coequalizer \(\operatorname{Coeg}(\varphi,\psi)\) and let \(\gamma\colon G\to G/K\) be the canonical groupoid morphism.

Let o be a vertex of G. For \(k \in K\), choose an arrow \(c_{k}\) connecting \(k^{*}o\) to o in G; such an arrow exists because G is transitive. For \(k \in K\), denote by \(\text{Fix}(k)\) the subgroupoid of G whose vertices (resp. arrows) are the elements of \(\text{Som}(G)\) (resp. of \(\text{Fl}(G)\) ) fixed by k; let us choose a set \(A_{k}\) of vertices of \(\text{Fix}(k)\) meeting all the orbits of this groupoid. For \(k \in K\) and \(a \in A_{k}\) we also choose an arrow \(f_{(a,k)}\) in G connecting the vertex a to the vertex o.

Proposition 6. — The unique group homomorphism

\[
\lambda \colon \mathrm{G} _ {o} * \mathrm{F} (\mathrm{K}) \to (\mathrm{G} / \mathrm{K}) _ {\gamma (o)}
\]

such that \(\lambda(f) = \gamma_o(f)\) for \(f \in \mathrm{G}_o\) and \(\lambda([k]) = \gamma(c_k)\) for \(k \in \mathrm{K}\) is surjective. Its kernel is the smallest normal subgroup of \(\mathrm{G}_o * \mathrm{F}(\mathrm{K})\) containing the following elements:

\[
r _ {2} (k, f) = [ k ] ^ {- 1} f [ k ] (c _ {k} ^ {- 1} k ^ {*} (f) ^ {- 1} c _ {k})\tag{\((R_{2})\)}
\]

\[
\text { for } k \in \mathrm{K} \text { and } f \in \mathrm{G} _ {o};\tag{\((\mathrm{R}_{3}^{\prime})\)}
\]

\[
r _ {3} ^ {\prime} (k, a) = [ k ] (c _ {k} ^ {- 1} k ^ {*} (f _ {(a, k)}) ^ {- 1} f _ {(a, k)}))
\]

\[
\text { for } k \in \mathrm{K} - \{e \} \text { and } a \in \mathrm{A} _{k}\tag{\((\mathrm{R}_{3}^{\prime\prime})\)}
\]

\[
r _ {3} ^ {\prime \prime} (k, h) = [ k h ] ^ {- 1} [ k ] [ h ] (c _ {h} ^ {- 1} h ^ {*} (c _ {k} ^ {- 1}) c _ {k h})
\]

\[
\text { for } k,h \in \mathrm{K}.
\]

Since G is transitive, the map which associates to an element \(k \in K\) the orbit of \((o, k)\) is a bijection from K onto the set of orbits of H. We thus identify \(\operatorname{Orb}(\mathbf{G})\) with \(\{o\}\) and \(\operatorname{Orb}(\mathbf{H})\) with K. The frame \(\Gamma_{0}\) of the pair \((\varphi, \psi)\) then identifies with the quiver having a single vertex o and whose set of arrows is K. Let T be the subquiver of \(\Gamma_{0}\) of vertex set \(\{o\}\) and whose set of arrows is empty; the associated graph is the unique maximal tree of the graph \(\widetilde{\Gamma}_{0}\).

The family \((o,(o,k)_{k\in \mathrm{K}},(e_o)_{k\in \mathrm{K}},(c_k)_{k\in \mathrm{K}},\mathrm{T},o)\) is a base equipment of the pair \((\varphi ,\psi)\) .

The map \(f \mapsto (f, k)\) defines an isomorphism of the isotropy group \(G_o\) onto the isotropy group \(H_{(o,k)}\) ; under this isomorphism, the homomorphisms \(\varphi_k\) and \(\psi_k\) of \(H_{(o,k)}\) into \(G_o\) defined by prop. 5 of II, p. 208 are given by

\[
\varphi_ {k} (f, k) = f \quad \text { and } \quad \psi_ {k} (f, k) = c _ {k} ^ {- 1} k ^ {*} (f) c _ {k},\tag{3}
\]

for \(k\in \mathrm{K}\) and \(f\in \mathrm{G}_o\).

The quiver \(H \times_{G} H\) is the equalizer in \(H \times H\) of quiver morphisms \(\psi \circ pr_{1}\) and \(\varphi \circ pr_{2}\) . Its vertices are the pairs \(((a, k), (b, h))\) where a and b are vertices of G and k and h are elements of K such that \(b = k^{*}a\) , and its arrows are the pairs \(((f, k), (g, h))\) where f and g are arrows of G, and k and h are elements of K such that \(g = k^{*}f\) ; the origin of the arrow \(((f,k),(g,h))\) is equal to \(((o(f),k),(o(g),h))\) ; its target is equal to \(((t(f),k),(t(g),h))\) .

We then define a quiver morphism \(\mu\) of \(H \times_{\mathrm{G}} H\) into H by setting \(\mu((a,k),(b,h)) = (a,kh)\) and \(\mu((f,k),(g,h)) = (f,kh)\) . We have \(\varphi \circ \mu = \varphi \circ \mathrm{pr}_1\) and \(\psi \circ \mu = \psi \circ \mathrm{pr}_2\) .

Let \(x\) and \(y\) be vertices of H such that \(\varphi(x) = \varphi(y)\). There thus exists a vertex \(a\) of G and elements \(k\) and \(h\) of K such that \(x = (a, k)\) and \(y = (a, h)\) . Set \(z = (h^* a, h^{-1} k)\) ; we have \(\mu(y, z) = x\). This shows that the pair \((\varphi, \psi)\) satisfies the hypotheses of Prop. 4 (II, p. 202).

A vertex \((a,k)\) of H belongs to the groupoid \(\operatorname{Ker}(\varphi,\psi)\), equalizer of \(\varphi\) and \(\psi\), if and only if \(k^{*}a=a\), that is, if k belongs to the stabilizer of the vertex a in the group K. Let \(A_{1}\) be the set of pairs \((a,k)\), for \(k\in K\) and \(a\in A_{k}\); it meets all the orbits of the subgroupoid \(\operatorname{Ker}(\varphi,\psi)\) of H. Let \(Z_{1}\) be the part of \(\Omega(\widetilde{\Gamma})\) formed of sequences of the form (((a,k),1)), for \((a,k)\in A_{1}\). For \(\mathbf{z}=((a,k),1)\in\mathbf{Z}_{1}\), \((f_{(a,k)},k)\) is an arrow of H which connects \((a,k)\) to \((o,k)\). Set \(f(\mathbf{z})=((f_{(a,k)},k))\).

The set \(\mathrm{A}_2\) of vertices of \(\mathrm{H} \times_{\mathrm{G}} \mathrm{H}\) of the form \(((o,k), (k^*o,h))\) , for \((k,h) \in \mathrm{K}^2\) , meets every orbit of \(\mathrm{H} \times_{\mathrm{G}} \mathrm{H}\) . Observe also that one has \(\mu((o,k), (k^*o,h)) = (o,kh)\) . The arrows \((e_o,k), (c_k,h), (e_o,kh)\) in H connect respectively \((o,k), (k^*o,h), (o,kh)\) to \((o,k), (o,h)\) and \((o,kh)\) . Denote by \(\mathrm{Z}_2\) the set of sequences of the form \((((o,k),1), ((k^*o,h), 1), ((o,kh), -1))\) in \(\Omega(\widetilde{\Gamma})\) ; for such an element \(\mathbf{z}\) of \(\mathrm{Z}_2\) , set \(f(\mathbf{z}) = ((e_o,k), (c_k,h), (e_o,kh))\) .

By prop. 4 of II, p. 202, the set \(Z = Z_1 \cup Z_2\) and the family \((f(\mathbf{z}))_{z \in \mathbb{Z}}\) form a complementary equipment.

Let \(k \in \mathrm{K}\) and let \(a \in \mathrm{A}_k\) . The element \(\mathrm{C}((a, k), 1)\) of the group \(\mathrm{G}_o * \mathrm{F}(\mathrm{K})\) defined by formula (2) of II, p. 208 is equal to

\[
c _ {k} ^ {- 1} \psi ((f _ {(a, k)}, k)) ^ {- 1} \varphi (f _ {(a, k)}) e _ {o} [ k ] = (c _ {k} ^ {- 1} k ^ {*} (f _ {(a, k)}) ^ {- 1} f _ {(a, k)}) [ k ].\tag{4}
\]

Let \(k,h \in \mathbf{K}\). One verifies that the element

\[
\mathrm{C} (((o, k), 1), ((k ^ {*} o, h), 1), ((o, k h), - 1))
\]

of the group \(G_{o} * F(K)\) defined by formula (2) of II, p. 208 is equal to

\[
[ k ] [ h ] (c _ {h} ^ {- 1} h ^ {*} (c _ {k} ^ {- 1}) c _ {k h}) [ k h ] ^ {- 1}.\tag{5}
\]

The elements \(r_3'(e, a)\) given by the relations \((\mathrm{R}_3')\) for \(k = e\) are all equal to \([e]c_e^{-1}\), an element of the group \(\mathrm{G}_o * \mathrm{F}(\mathrm{K})\) obtained by applying the relation \((\mathrm{R}_{3}^{\prime\prime})\) to \(k = h = e\). In view of relations (3), (4), and (5), the proposition therefore follows from II, p. 208, prop. 5.

COROLLARY 1. --- Suppose that the group K is generated by the fixators of the vertices of G. Then the group homomorphism \(\gamma_{o}\colon\mathrm{G}_{o}\to(\mathrm{G}/\mathrm{K})_{\gamma(o)}\) is surjective. Moreover, if the groupoid G is simply transitive, then so is the groupoid G/K.

The relation \((\mathrm{R}_{3}^{\prime\prime})\) implies \(\lambda([e]) = \lambda(c_{e}) = \gamma_{o}(c_{e})\) . The relations \((\mathrm{R}_{3}^{\prime\prime})\) then imply that the set of \(k \in K\) such that \(\lambda([k])\) belongs to the image of \(\gamma_{o}\) is a subgroup of K. Finally, the relations \((\mathrm{R}_{3}^{\prime})\) show that, for every element \(k \in K\) whose fixator is not empty, \(\lambda([k])\) belongs to the image of \(\gamma_{o}\) . The corollary follows from this.

Remark 1. --- One can give another description, sometimes more convenient, of the group \((\mathrm{G}/\mathrm{K})_{\gamma(o)}\). For this, set \(M = K \times G_{o}\) and let us define a composition law on M by the formula

\[
(k, a) \cdot (h, b) = (k h, c _ {k h} ^ {- 1} h ^ {*} (c _ {k} a) c _ {h} b),
\]

for \(k\) , \(h \in \mathrm{K}\) and \(a\) , \(b \in \mathrm{G}_o\) . We verify that this composition law is associative, that \((e, c_e^{-1})\) is an identity element and that the element \((k^{-1}, c_{k^{-1}}^{-1}(k^{-1})^*(c_k a)^{-1} c_e)\) is the inverse of \((k, a)\) . It therefore endows M with a group structure. Moreover, the map \(\lambda'\colon \mathrm{M} \to (\mathrm{G} / \mathrm{K})_{\gamma(o)}\) defined by \((k, a) \mapsto \lambda([k]a)\) is a group homomorphism. Let \(\alpha'\) be the unique group homomorphism from \(\mathrm{G}_o * \mathrm{F}(\mathrm{K})\) to M such that \(\alpha'(f) = (e, f)\) if \(f \in \mathrm{G}_o\) and \(\alpha'([k]) = (k, e_o)\) if \(k \in \mathrm{K}\) ; we have \(\lambda' \circ \alpha' = \lambda\) .

The relations \((\mathrm{R}_{2})\) and \((\mathrm{R}_{3}^{\prime\prime})\) show that every element of \((\mathrm{G}/\mathrm{K})_{\gamma(o)}\) is the image under the homomorphism \(\lambda\) of an element of \(\mathrm{G}_{o}*\mathrm{F}(\mathrm{K})\) of the form \([k]f\) , with \(f\in G_{o}\) and \(k\in K\) . Consequently, the homomorphism \(\lambda'\) is surjective. One also verifies that the image under \(\alpha'\) of an element of \(\mathrm{G}_{o}*\mathrm{F}(\mathrm{K})\) of the form \((\mathrm{R}_{2})\) or \((\mathrm{R}_{3}^{\prime\prime})\) is zero. Consequently, the kernel of the homomorphism \(\lambda'\) is the smallest normal subgroup of M containing the images under \(\alpha'\) of the elements of \(\mathrm{G}_{o}*\mathrm{F}(\mathrm{K})\) of the form \((\mathrm{R}_{3}^{\prime})\) .

COROLLARY 2. --- Suppose that the group K acts freely on Som(G). Then there exists a unique group homomorphism \(\pi\colon (\mathrm{G} / \mathrm{K})_{\gamma (o)}\to \mathrm{K}\) whose kernel contains the image of \(\gamma_{o}\) and such that \(\pi (\lambda ([k])) = k\) for all \(k\in \mathbf{K}\) . Moreover, \(\mathrm{G}_o\xrightarrow{\gamma_o} (\mathrm{G} / \mathrm{K})_{\gamma (o)}\xrightarrow{\pi}\mathrm{K}\) is an extension of K by \(\mathrm{G}_o\) .

If such a group homomorphism \(\pi\) exists, the group homomorphism \(\pi \circ \lambda\) is necessarily equal to the unique group homomorphism \(p\) of \(\mathrm{G}_{o} * \mathrm{F}(\mathrm{K})\) into \(\mathrm{K}\) such that \(p(f) = e\) for \(f \in \mathrm{G}_{o}\) and \(p([k]) = k\) . It is immediate to verify that the elements of \(\mathrm{G}_{o} * \mathrm{F}(\mathrm{K})\) defined by the formulas \((\mathrm{R}_{2})\) and \((\mathrm{R}_{3}^{\prime \prime})\) belong to the kernel of \(p\) . By hypothesis, there is no element of the type \((\mathrm{R}_{3}^{\prime})\) . Thus, the kernel of the morphism \(\lambda\) contains that of \(p\) . Consequently, there exists a unique group homomorphism \(\pi: (\mathrm{G} / \mathrm{K})_{\gamma(o)} \to \mathrm{K}\) such that \(\pi \circ \lambda = p\) .

It is evident that the homomorphism \(\pi\) is surjective. To show that the homomorphism \(\gamma_{o}\) is injective and that its image is exactly the kernel of \(\pi\), let us note that the homomorphism \(\lambda'\colon M \to (\mathrm{G}/\mathrm{K})_{\gamma(o)}\) (II, p. 213, remark 1) is an isomorphism, since we have assumed that K acts freely in \(\operatorname{Som}(\mathrm{G})\). The composite homomorphism \((\lambda')^{-1} \circ \gamma_{o}\) of \(\mathrm{G}_{o}\) into M is given by \(f \mapsto (e, f)\), while the homomorphism \(\pi \circ \lambda'\colon M \to K\) maps \((k, f)\) to \(k\). The corollary follows from this.

COROLLARY 3. --- Suppose that the groupoid G is simply transitive. Let K₀ be the subgroup of K generated by the stabilizers of the vertices of G. The map from K to \((\mathrm{G}/\mathrm{K})_{\gamma(o)}\) which sends \(k\) to \(\gamma(c_k)\) is a surjective group homomorphism, with kernel K₀.

If an element \(k \in \mathrm{K}\) fixes a vertex \(a\) of \(\mathrm{G}\) , the element \(g^{-1}kg\) fixes the vertex \(g^{*}a\) ; this implies that \(\mathrm{K}_{0}\) is a normal subgroup of \(\mathrm{K}\) .

By proposition 5, the unique homomorphism \(\lambda\colon\mathrm{F}(\mathrm{K})\to(\mathrm{G}/\mathrm{K})_{\gamma(o)}\) such that \(\lambda([k])=\gamma(c_k)\) is surjective, and its kernel is the smallest normal subgroup of F(K) containing the elements \([k]\), where \(k\) is an element of K that fixes a vertex of G, and the elements \([kh]^{-1}[k][h]\) , where \((k,h)\in\mathrm{K}^2\) . In particular, the map \(\lambda'\colon\mathrm{K}\to(\mathrm{G}/\mathrm{K})_{\gamma(o)}\) , defined by \(\lambda'(k)=\gamma(c_k)=\lambda([k])\) for \(k\in\mathrm{K}\) , is a group homomorphism. We have \(\lambda=\lambda'\circ p\) , where \(p\colon\mathrm{F}(\mathrm{K})\to\mathrm{K}\) denotes the canonical surjective group homomorphism. Consequently, the homomorphism \(\lambda'\) is surjective and its kernel is the smallest normal subgroup of K that contains the elements \(k\) , for \(k\) fixing a vertex of G, that is, \(\mathrm{K}_0\). This proves the proposition.
% semantic-fragments: legacy-input-seam
\relax{}
